\chapter{Baire spaces}\label{chap:baire}

Here we will develop a language to talk rigorously about typical objects in many contexts.
In practice, it is used to prove existence of examples with weird properties,
and it often comes with a surprise: typical object looks like nothing we saw before.

\section{Baire theorem}

A topological space is called \index{Baire space}\emph{Baire} if any countable family of dense open sets in $\c{X}$ has a dense intersection.

For example, the space of rational numbers $\QQ$ is not Baire.
Indeed, the complement of a one-point set $\QQ\setminus \{q\}$ is dense and open in $\QQ$;
the intersection of all these complements is empty, and since $\QQ$ is countable, this is a countable intersection.
On the other hand the set of integers $\ZZ$ is a Baire space, since $\ZZ$ is the only dense open set in $\ZZ$.
The following theorem implies that $\RR$ is a Baire space.


\begin{thm}{Theorem}\label{thm:baire-category}
Any complete metric space is Baire.
\end{thm}

Given two subsets $A$ and $B$ of a topological space, we will write \index{$A\Supset B$}$A\Supset B$ (equivalently, $B\Subset A$) if $\mathring{A}\supset \bar B$; that is, the interior of $A$ contains the closure of $B$.

\parit{Proof.}
Let $V_1,V_2\ldots$ be a sequence of dense open sets in a complete metric space $\c{M}$.
We need to show that $G=V_1\cap V_2\cap\ldots$ intersects any open ball in $\c{M}$.

Choose a ball $B_0=\Ball(p_0,r_0)$.
Since $V_1$ is dense and open, we can choose another ball $B_1=\Ball(p_1,r_1)\subset B_0\cap V_1$; moreover, taking the radius $r_1$ smaller we can achieve
$B_1\Subset B_0\cap V_1$ and $r_1\le\tfrac12\cdot r_0$.
Continue this process for $V_2$, $V_3$, and so on;
on the $n$-th step we choose a ball $B_n=\Ball(p_n,r_n)\Subset B_{n-1}\cap V_n$ and $r_n\le\tfrac12\cdot r_{n-1}$.

If $i<j$, then  $|p_i-p_j|<r_i\le\tfrac 1{2^i}\cdot r_0$.
It follows that $p_0,p_1,\ldots$ is a Cauchy sequence.
Since $\c{M}$ is complete, the sequence must converge; denote its limit by $p_\infty$.
It remains to observe that $p_\infty\in \bar B_n$ for any $n$; in particular, $p_\infty\in G\cap B_0$.
\qeds


\begin{thm}{Exercise}\label{ex:haus-comp-bair}
Let $\c{K}$ be a compact Hausdorff space.
\begin{subthm}{ex:haus-comp-bair:Subset}
Given a nonempty open set $V\subset \c{K}$, construct a nonempty open set $W\Subset V$.
\end{subthm}
\begin{subthm}{ex:haus-comp-bair:bair}
Conclude that $\c{K}$ is a Baire space.
\end{subthm}

\end{thm}

\begin{thm}{Exercise}\label{ex:irrational-baire}
Show that irrational numbers $\RR\setminus\QQ$ form a Baire space.
\end{thm}

\begin{thm}{Exercise}\label{ex:union-closures}
Let $B_1,B_2,\dots$ be a sequence of nonempty sets in a Baire space $\c{X}$.
Assume $\c{X}=B_1\z\cup B_2\cup\dots$
Show that at least one of the closures $\bar B_n$ has nonempty interior.
\end{thm}

\begin{thm}{Exercise}\label{ex:Baire-subsets}
Prove the following statements.

\begin{subthm}{ex:Baire-subsets:open}
Any open set in Baire space is Baire.
\end{subthm}

\begin{subthm}{ex:Baire-subsets:loc}
Any locally Baire space is Baire; that is, if any point $p$ in a topological space $\c{X}$ has a Baire neighborhood, then $\c{X}$ is Baire.
\end{subthm}

\end{thm}


\section{G-delta sets}

Let $G$ be a set in a topological space $\c{X}$.
We say that $G$ is a \index{G-delta set}\emph{G-delta} set%
if it can be written as a countable intersection of open sets;
that is, $G=V_1\cap V_2\cap\ldots$ for a sequence of open sets $V_1,V_2,\ldots \subset \c{X}$.%
In the definition we can assume that the sequence of open sets is nested;
that is $V_1\supset V_2\supset\ldots$;
if not, pass to the sequence
\[V_1,\ V_1\cap V_2,\ V_1\cap V_2\cap V_3,\ \dots\]

\begin{thm}{Exercise}\label{ex:closed-G-delta}
Show that any closed set in a metrizable space is G-delta.
\end{thm}

Observe that a countable intersection of G-delta sets is G-delta.
The following identity implies that  union of two (and therefore finitely many) G-delta sets is G-delta:
\[\biggl(\bigcap_i V_i\biggr)\cup \biggl(\bigcap_j W_j\biggr)
=
\bigcap_{i,j}\, (V_i\cup W_j).\]

\begin{thm}{Observation}\label{obs:G-delta}
Let $\c{X}$ be a Baire space.
Then countable intersection of dense G-delta sets in $\c{X}$ is dense.
\end{thm}

\begin{thm}{Exercise}\label{ex:G-delta=closed n dense-G-delta}
\begin{subthm}{ex:G-delta=closed n dense-G-delta:closed}
Show that any G-delta set is an intersection of a closed set with a dense G-delta set.
\end{subthm}

\begin{subthm}{ex:G-delta=closed n dense-G-delta:open}
Give an example of a topological space with a G-delta set that is not an intersection of an open set with a dense G-delta set.
\end{subthm}

\end{thm}

Note that the set of irrational numbers $\II=\RR\setminus \QQ$ is G-delta --- it is an intersection of all complements $\RR\setminus\{q\}$ for $q\in\QQ$.
Since $\QQ$ and $\II$ are dense, the observation implies that $\QQ\subset\RR$ is not G-delta.
On the other hand,  $\QQ\subset\RR$ is a countable union of one-point sets; each one-point set $\{q\}$ can be written as intersection of open intervals $(q-\tfrac1n,q+\tfrac1n)$.
Therefore, countable union of G-delta sets does not have to be G-delta.

\begin{thm}{Exercise}\label{ex:uncountable-Gdelta}
Show that any dense G-delta set in $\RR$ is uncountable.
\end{thm}

Let us consider a map $f\:\c{X}\to\c{Y}$ between topological spaces.
We say that $f$ is \index{continuous at a point}\emph{continuous} at $x\in \c{X}$ if for any open set $W\subset \c{Y}$ such that $W\ni f(x)$ there is an open set $V\ni x$ such that $f(V)\subset W$.
Observe that $f$ is continuous if and only if it is continuous at every point $x\in \c{X}$.

\begin{thm}{Exercise}\label{ex:continuity-Gdelta}
Given a function $f\: \RR\to \RR$, denote by $\Cont f\subset \RR$ the set of points at which $f$ is continuous.

\begin{subthm}{ex:continuity-Gdelta:cont}
Show that $\Cont f$ is a G-delta set for any function $f\: \RR\to \RR$.
\end{subthm}

\begin{subthm}{ex:continuity-Gdelta:R-Q}
Show that there is no function $f\: \RR\to \RR$ such that $\Cont f=\QQ$.
\end{subthm}

\begin{subthm}{ex:continuity-Gdelta:Q}
Let us enumerate all rational numbers $r_1,r_2,\dots$
Consider function $f\: \RR\to \RR$ defined by $f(r_n)=\tfrac1n$ and $f(x)=0$ for any irrational $x$.
Show that $\Cont f=\RR\setminus\QQ$.
\end{subthm}

\end{thm}

\begin{thm}{Advanced exercise}\label{ex:Baire-closed-subsets}
Show that any G-delta set in a complete metric space is Baire.

Give an example of a Baire space $\c{X}$ with a G-delta set that is not Baire.
\end{thm}






\section{Large or small?}

According to the observation, dense G-delta sets contain all points in Baire space except a \textit{negligible} set.
At least, we cannot cover an open set with complements of dense G-delta sets
and intersection of countably many dense G-delta sets is a dense G-delta set.

In spaces with measure, the sets of full measure have analogous features, but as we will see in an incompatible way.

Let us enumerate rational numbers $q_1,q_2,\dots$, consider the unions of open intervals:
\[V_n\df \bigcup_i (q_i-\tfrac1{2^{n+i}},q_i+\tfrac1{2^{n+i}}).\]
Note that $G=V_1\cap V_2\cap \ldots$ is a dense G-delta set in $\RR$.
Observe that the Lebesgue measure of $V_n$ (its total length) is less than
\[\frac1{2^{n-1}}=\frac{2}{2^{n+1}}+\frac{2}{2^{n+2}}+\ldots\]
Therefore, $G$ has vanishing measure.
So, measure-theoretically, $G$ is nearly nothing, but topologically it is nearly everything.

\section{An application}

Let $\c{C}$ be the set of all continuous functions $[0,1]\to\RR$,
equipped with the \index{sup-norm metric}\emph{sup-norm metric}
\[|f-g|_{\c{C}}\df\sup\set{|f(x)-g(x)|}{x\in [0,1]}.\]
It is straightforward to check that this is indeed a metric;
that is, it meets all conditions in \ref{def:metric-space}.
The following claim is not that obvious.

\begin{thm}{Claim}
The space $\c{C}$ is complete.
\end{thm}

\parit{Proof.}
In the space $\c{C}$ every point is a continuous function $[0,1]\to\RR$.
We need to show that a Cauchy sequence $f_1,f_2,\ldots \in \c{C}$ converges in~$\c{C}$.

Choose $x\in [0,1]$.
Observe that $f_1(x), f_2(x),\dots$ is a Cauchy sequence of real numbers; so,
it must converge; let us denote its limit by $f_\infty(x)$.

Evidently
\[\sup\set{|f_n(x)-f_\infty(x)|}{x\in[0,1]}\to 0.\eqlbl{eq:|f-f|->0}\]
Fix $x_0\in [0,1]$ and $\eps>0$.
By \ref{eq:|f-f|->0}, we can choose $n$ such that
\[|f_n(x)-f_\infty(x)|<\tfrac\eps3\]
for any $x\in [0,1]$.
By continuity of $f_n$, we can choose $\delta>0$ such that
\[|x-x_0|<\delta
\quad\Rightarrow\quad |f_n(x)-f_n(x_0)|<\tfrac\eps3\]
If $|x-x_0|<\delta$, then
\begin{align*}
\eps&>|f_\infty(x)-f_n(x)|+|f_n(x)-f_n(x_0)|
+|f_n(x_0)-f_\infty(x_0)|\ge
\\
&\ge
|f_\infty(x)-f_\infty(x_0)|.
\end{align*}
Since $\eps>0$ and $x_0$ are arbitrary, $f_\infty$ is continuous; that is, $f_\infty\in\c{C}$.
Finally by \ref{eq:|f-f|->0}, $f_n\to f_\infty$ in $\c{C}$.
\qeds

\begin{thm}{Claim}
The set of all continuous functions $[0,1]\to\RR$ that are not differentiable at any point contains a dense G-delta set in $\c{C}$.
\end{thm}

\parit{Sketch of proof.}
Given a positive integer $n$ consider the set $V_n$ of all functions $f$ in $\c{C}$ such that
for some integer $m>n$ we have the following inequality
\[|f(\tfrac{k}{m})-f(\tfrac{k-1}{m})|>\tfrac nm\]
for any integer $k$ from $1$ to $m$.

The set $V_n$ is open and dense in $\c{C}$;
to prove the latter statement, take huge value $m$ and approximate a given function by zig-zag-looking function with large slope as shown in the picture.%
\footnote{More formally, let us choose small $\eps>0$.
Given a continuous function $f$ we can choose large $m$ so that if $|x_0-x_1|\le \tfrac1m$, then $|f(x_0)-f(x_1)|<\eps$.
Now define a new function $h$ by setting $h(\tfrac im)=f(\tfrac im)+(-1)^i\cdot\eps$ for integer $i$ from $0$ to $m$ and extend $h$ linearly on each interval $[\tfrac {i-1}m,\tfrac im]$.
Observe that $h\in \c{C}$ and $|h-f|_{\c{C}}<2\cdot \eps$; moreover given $\eps>0$ and a positive integer $n$, one can choose $m$ sufficiently large so that $h\in V_n$.}

{

\begin{wrapfigure}{o}{26 mm}
\vskip-0mm
\centering
\includegraphics{mppics/pic-940}
\end{wrapfigure}

Now, assume $f\in V_n$ for every $n$.
Observe that for any $x_0\in [0,1]$ we can choose a point $x\in [0,1]$ such that
\[|x-x_0|< \tfrac1n,
\quad\text{and}\quad
\left|\frac{f(x)-f(x_0)}{x-x_0}\right|>n.
\]
It follows that the limit
\[\lim_{x\to x_0}\frac{f(x)-f(x_0)}{x-x_0}\]
is undefined.
Therefore $f$ is not differentiable at every $x_0\in [0,1]$.

}

It remains to apply \ref{thm:baire-category} to $G=V_1\cap V_2\cap\ldots$
\qeds

We just proved that \textit{most} continuous functions $[0,1]\z\to\RR$ are not differentiable at any point.
Try to come up with an explicit example of one such function --- it is not easy.
Maybe you did not see such functions before, altho you know many functions; in this case the functions you know are very non-typical and you did not meet a typical one.

Let me list several more statements that can be proved essentially the same way; try to think of explicit examples of that type.
\begin{itemize}
\item A typical surjective nondecreasing function $f\:[0,1]\to[0,1]$ is strictly increasing, but its derivative vanishes at every point where it is defined.
\item  A typical distance-nonexpanding map $f\:\SS^2\z\to\RR^2$ is \index{length-preserving map}\emph{length-preserving}; that is,
\[\length f\circ\gamma=\length \gamma\]
for any curve $\gamma\:[a,b]\to \SS^2$.
\item  A typical compact set in the plane is homeomorphic to the Cantor set; see \ref{sec:Cantor set}. Furthermore, a typical compact \textit{connected} set (defined in the next chapter) is homeomorphic to the pseudo-arc, which is a truly weird example in many ways (read about~it).
\end{itemize}

%the uniform boundedness principle

\section{Borel sets}

Any topological space $\c{X}$ comes with the so-called \index{Borel sets}\emph{Borel sets}, which form the minimal family of sets that is closed under complement and countable union and includes all open sets in $\c{X}$.

Borel sets play a key role in analysis.
We will discuss Borel sets in this section, but we are not going to use them further in the text.

A family of sets in $\c{X}$ that contains the whole space and is closed under complements and countable unions is called a \index{sigma-algebra}\emph{sigma-algebra}, so we can say that Borel sets form the minimal sigma-algebra that includes all open sets.

Since complement of countable union is the countable intersection of the complements,
sigma-algebra can be equivalently defined using countable intersection instead of countable union.
It follows that Borel sets include all closed sets (since they are complements of open sets) as well as G-delta sets (as countable intersections of open sets) and their complement, which are called \index{F-sigma set}\emph{F-sigma} sets;
the F-sigma sets can be defined directly as union of countable families of closed sets.


The following proposition says that any Borel set in $\RR$ is a negligible modification of an open set.

\begin{thm}{Proposition}\label{prop:B=GnV}
For any Borel set $B$ in a Baire space there is a dense G-delta set $G$ and an open set $V$ such that $B\cap G=V\cap G$.
\end{thm}

\parit{Proof.}
Choose a Baire space $\c{X}$.
Let us denote by $\s{B}$ the set of all Borel sets in $\c{X}$.
Further, denote by $\s{S}$ the set of all subsets of $\c{X}$ that meet the conclusion of the proposition;
that is, a set $S\subset \c{X}$ belongs to $\s{S}$ if there is a dense G-delta set $G\subset\c{X} $ and an open set $V\subset\c{X}$ such that $S\cap G=V\cap G$.

Observe that $\s{S}$ includes all open sets.
Let us show that $\s{S}$ is a sigma-algebra;
once it is proved, the proposition follows since $\s{B}$ is the minimal sigma-algebra that includes all open sets in $\c{X}$.

Choose $S\in \s{S}$, so there is a dense G-delta set $G$ and an open set $V$ such that $S\cap G=V\cap G$.
Let $T=\c{X}\setminus S$.
Note that $W=\c{X}\setminus \bar V$ is open, $\c{X}\setminus\partial V=V\cup W$ is a dense open set, and therefore $G'=G\cap (V\cup W)$ is a dense G-delta set.
It follows that
\[T\cap G'=W\cap G',\]
so $T\in \s{S}$, and we proved that $\s{S}$ is closed under taking complements.

Now, choose a sequence of sets $S_1,S_2,\ldots\in \s{S}$; it comes with a sequence of dense G-delta sets $G_1,G_2,\ldots$ and open sets $V_1,V_2,\dots$ such that $S_i\cap G_i=V_i\cap G_i$ for each $i$.
Let
\[
S=S_1\cup S_2\cup \ldots,\quad
V=V_1\cup V_2\cup \ldots,\quad
G=G_1\cap G_2\cap \ldots
\]
The set $V$ is open as a union of open sets.
Furthermore, $G$ is a G-delta set, and since $\c{X}$ is Baire, it must be dense.
Observe that
\[S\cap G=V\cap G;\]
in particular, $S\in \s{S}$ --- we proved that $\s{S}$ is closed under countable union, so $\s{S}$ is a sigma-algebra --- the proof is finished.
\qeds

We say that $L\subset\RR$ is a \index{Vitali set}\emph{Vitali set} if for any $x\in\RR$ there is a unique $\ell\in L$ such that $x-\ell\in \QQ$; here $\QQ$ denotes the set of rational numbers.

To construct a Vitali set, consider equivalence relation $\sim$ on $\RR$ such that $x\sim y$ $\Leftrightarrow$ $x-y\in \QQ$
and choose one number from each equivalence class.
However, to \textit{choose} a number in every equivalence class one needs the axiom of choice.

Now we are ready to show that not every set in the real line is Borel.

\begin{thm}{Claim}\label{clm:vitali-borel}
A Vitali set $L\subset\RR$ is not Borel.
\end{thm}

\parit{Proof.}
Given a set $S\subset \RR$, denote by $x+S$ its shift by $x\in\RR$; that is,
\[x+S=\set{x+s}{s\in S}.\]

Let $L$ be a Vitali set.
Note that $L\cap (q+L)=\emptyset$
for any rational $q\ne 0$
and
\[\RR=\bigcup_{q\in \QQ}(q+L).\eqlbl{eq:cupL}\]

If $L$ is Borel, then by \ref{prop:B=GnV} $L\cap G=V\cap G$ for some dense G-delta set $G$ and an open set $V$.

Note that $V\ne \emptyset$;
otherwise \ref{eq:cupL} implies that
\[\bigcap_{q\in \QQ}(q+G)\subset \RR\setminus \bigcup_{q\in Q}(q+L)=\emptyset,\]
which is impossible since $q+G$ is a dense G-delta set for any $q\in \QQ$, the set $\QQ$ is countable, and $\RR$ is Baire.

Since $V$ is open, it contains an open interval;
therefore, we can choose small $q\in\QQ$ such that $V_q=V\cap (q+V)$ is a nonempty open set.
Let $G_q=G\cap (q+G)$; it is a dense G-delta set.
Hence $V_q\cap G_q\ne\emptyset$.
Therefore
\[\emptyset=L\cap (q+L)\supset V_q\cap G_q\ne\emptyset \]
--- a contradiction.
\qeds

\begin{thm}{Exercise}\label{ex:nonborel-image}
Construct a continuous map $f\:\c{X}\to \RR$ and a Borel set $B\subset \c{X}$ such that $f(B)$ is not Borel.
\end{thm}

\begin{thm}{Advanced exercise}\label{ex:borel-square}
Construct a bijection $f\:[0,1]\to[0,1]\z\times [0,1]$ such that $f$ and its inverse $f^{-1}$ map Borel sets to Borel sets.
\end{thm}
